\documentclass{article}
\usepackage{geometry}
\geometry{margin=1.5cm, vmargin={0pt,1cm}}
\setlength{\topmargin}{-1cm}
\setlength{\headheight}{12.5pt}
\setlength{\paperheight}{29.7cm}
\setlength{\textheight}{25.3cm}

% useful packages.
\usepackage[UTF8]{ctex}
\usepackage{fancyhdr}
\usepackage{layout}
\usepackage{tikz}

% import common macros
% % % % % % % % % % % % % % % % % % % % % % % % % % % % % % % %
% Common Macros for LaTeX
% % % % % % % % % % % % % % % % % % % % % % % % % % % % % % % %

% Useful packages
\usepackage{amsfonts}
\usepackage{amsmath}
\usepackage{amssymb}
\usepackage{amsthm}
\usepackage{enumerate}
\usepackage{graphicx}
\usepackage{multicol}
\usepackage[ruled,linesnumbered]{algorithm2e}

% Math operators and common symbols
\newcommand{\dif}{\mathrm{d}}
\newcommand{\innerProd}[1]{\left\langle #1 \right\rangle}
\newcommand{\difFrac}[2]{\frac{\dif #1}{\dif #2}}
\newcommand{\pdfFrac}[2]{\frac{\partial #1}{\partial #2}}
\newcommand{\T}{\mathrm{T}}
\newcommand{\norm}[1]{\left\| #1 \right\|}
\newcommand{\abs}[1]{\left| #1 \right|}

% Vector and Matrix shortcuts
\newcommand{\vecTwo}[2]{
  \begin{bmatrix}
    #1 \\
    #2
\end{bmatrix}}
\newcommand{\vecThree}[3]{
  \begin{bmatrix}
    #1 \\
    #2 \\
    #3
\end{bmatrix}}
\newcommand{\vecTwoH}[2]{
  \begin{bmatrix}
    #1 & #2
  \end{bmatrix}
}
\newcommand{\vecThreeH}[3]{
  \begin{bmatrix}
    #1 & #2 & #3
  \end{bmatrix}
}

% Common fields
\newcommand{\R}{\mathbb{R}}
\newcommand{\C}{\mathbb{C}}
\newcommand{\Z}{\mathbb{Z}}
\newcommand{\N}{\mathbb{N}}

% Solutions and proofs
\newenvironment{solution}{
\begin{proof}[解]}{
\end{proof}}

% Utility to get environment variables (requires lualatex)
\newcommand{\getEnv}[2]{\directlua{tex.print(os.getenv("#1") or "#2")}}


% Name and uID
\newcommand{\name}{\getEnv{NAME}{NAME}}
\newcommand{\uid}{\getEnv{UID}{UID}}
\newcommand{\course}{Real Analysis}
\newcommand{\hwNum}{16}

\begin{document}

\pagestyle{fancy}
\fancyhead{}
\lhead{\zhstr{\name} (\uid)}
\chead{\course\ \#\hwNum}
\rhead{\today}

% ==============================================================================
% Page 135, Exercise 8
% ==============================================================================
\section{P136 Exercise 8}

设 $1<p<\infty$.
\begin{enumerate}
  \renewcommand{\labelenumi}{(\roman{enumi})}
  \item 证明 $\dfrac{\sin x}{x}\in L^p((0,\infty))$;
  \item 若 $f\in L^p((0,\infty))$, 求证:
    \[
      g(x)=\int_0^\infty f(t)\frac{\sin xt}{t}\,\dif t
    \]
    存在且 $h^{-1/p}[g(x+h)-g(x)]$ 有界, $x>0,\ h>0$.
\end{enumerate}

\begin{solution}
  (i)\; 在 $(0,1)$ 上, $\left|\dfrac{\sin x}{x}\right|\le1$.
  在 $(1,\infty)$ 上, $\left|\dfrac{\sin x}{x}\right|^p\le x^{-p}$.
  由于 $p>1$, $\int_1^\infty x^{-p}\,\dif x<\infty$, 故
  \[
    \frac{\sin x}{x}\in L^p((0,\infty)).
  \]

  (ii)\; 令 $q$ 为 $p$ 的共轭指数. 对任意固定的 $x>0$,
  \[
    \int_0^\infty\left|\frac{\sin xt}{t}\right|^q\,\dif t
    =x^{q-1}\int_0^\infty\left|\frac{\sin u}{u}\right|^q\,\dif u<\infty,
  \]
  其中最后一步由 (i) 对指数 $q>1$ 的情形得到. 因此
  $\dfrac{\sin xt}{t}\in L^q((0,\infty))$. 由 Hölder 不等式,
  $g(x)$ 对每个 $x>0$ 都存在.

  又
  \[
    g(x+h)-g(x)
    =\int_0^\infty f(t)
      \frac{\sin((x+h)t)-\sin xt}{t}\,\dif t.
  \]
  由
  \[
    |\sin((x+h)t)-\sin xt|
    =2\left|\cos\left(xt+\frac{ht}{2}\right)\sin\frac{ht}{2}\right|
    \le 2\left|\sin\frac{ht}{2}\right|,
  \]
  可得
  \begin{align*}
    \left\|\frac{\sin((x+h)t)-\sin xt}{t}\right\|_q^q
    &\le
    2^q\int_0^\infty\frac{|\sin(ht/2)|^q}{t^q}\,\dif t  \\
    &=2h^{q-1}\int_0^\infty\left|\frac{\sin u}{u}\right|^q\,\dif u.
  \end{align*}
  所以存在只依赖于 $p$ 的常数 $C_p$, 使得
  \[
    |g(x+h)-g(x)|
    \le C_p\|f\|_p h^{1/p}.
  \]
  于是
  \[
    h^{-1/p}|g(x+h)-g(x)|\le C_p\|f\|_p,
  \]
  即 $h^{-1/p}[g(x+h)-g(x)]$ 关于 $x>0,\ h>0$ 有界.
\end{solution}

% ==============================================================================
% Page 135, Exercise 10
% ==============================================================================
\section{P136 Exercise 10}

设 $f\in L^2([0,1])$, $F(x)=\displaystyle\int_0^x f(t)\,\dif t$.
求证:
\[
  \|F\|_2\le \frac1{\sqrt2}\|f\|_2.
\]

\begin{solution}
  对每个 $x\in[0,1]$, 由 Cauchy--Schwarz 不等式,
  \[
    |F(x)|^2
    =\left|\int_0^x f(t)\,\dif t\right|^2
    \le x\int_0^x|f(t)|^2\,\dif t.
  \]
  两边关于 $x$ 积分并交换积分次序:
  \begin{align*}
    \|F\|_2^2
    &\le \int_0^1x\int_0^x|f(t)|^2\,\dif t\,\dif x  \\
    &=\int_0^1|f(t)|^2\int_t^1x\,\dif x\,\dif t  \\
    &=\frac12\int_0^1(1-t^2)|f(t)|^2\,\dif t
    \le\frac12\|f\|_2^2.
  \end{align*}
  开方即得
  \[
    \|F\|_2\le\frac1{\sqrt2}\|f\|_2. \qed
  \]
\end{solution}

% ==============================================================================
% Page 135, Exercise 11
% ==============================================================================
\section{P136 Exercise 11}

设 $f\in L^p(\mathbf R)$, $1\le p<\infty$,
$\dfrac1p+\dfrac1q=1$, $F(x)=\displaystyle\int_0^x f(t)\,\dif t$.
求证: 对任何 $x$,
\[
  F(x+h)-F(x)=o(|h|^{1/q})\qquad (h\to0).
\]

\begin{solution}
  先设 $1<p<\infty$. 对任意固定的 $x\in\mathbf R$,
  \[
    F(x+h)-F(x)=\int_x^{x+h}f(t)\,\dif t.
  \]
  由 Hölder 不等式,
  \[
    |F(x+h)-F(x)|
    \le |h|^{1/q}
      \left(\int_{I_h}|f(t)|^p\,\dif t\right)^{1/p},
  \]
  其中 $I_h$ 是端点为 $x$ 与 $x+h$ 的区间. 因为
  $|f|^p\in L^1(\mathbf R)$, 积分的绝对连续性给出
  \[
    \int_{I_h}|f(t)|^p\,\dif t\longrightarrow0
    \qquad(h\to0).
  \]
  因而
  \[
    \frac{|F(x+h)-F(x)|}{|h|^{1/q}}\longrightarrow0.
  \]

  当 $p=1$ 时 $q=\infty$, 即 $1/q=0$. 此时
  \[
    |F(x+h)-F(x)|\le\int_{I_h}|f(t)|\,\dif t\longrightarrow0,
  \]
  同样得到结论. \qed
\end{solution}

% ==============================================================================
% Page 135, Exercise 13
% ==============================================================================
\section{P136 Exercise 13}

设 $1<p<\infty$, $f,f_k\in L^p(E)\ (k\ge1)$. 求证: 当
$f_k(x)\to f(x)$ a.e. 且 $\sup_k\|f_k\|_p<\infty$ 时, 对任何
$g\in L^q(E)$, $\dfrac1p+\dfrac1q=1$, 有
\[
  \int_E f_k(x)g(x)\,\dif x
  \longrightarrow
  \int_E f(x)g(x)\,\dif x.
\]

\begin{solution}
  由 Fatou 引理,
  \[
    \|f\|_p^p
    \le \liminf_{k\to\infty}\|f_k\|_p^p<\infty,
  \]
  故 $f\in L^p(E)$. 记
  \[
    M=\sup_k\|f_k-f\|_p
    \le \sup_k\|f_k\|_p+\|f\|_p<\infty.
  \]

  先证结论对有界且支集测度有限的函数 $\varphi$ 成立. 设
  $A=\{x:\varphi(x)\ne0\}$, $m(A)<\infty$. 由于 $f_k-f\to0$ a.e. 于
  $A$ 上, 由 Egorov 定理, 任给 $\delta>0$, 存在
  $B\subset A$ 且 $m(B)<\delta$, 使得 $f_k-f$ 在 $A\setminus B$ 上一致收敛到 $0$.
  因此
  \[
    \int_{A\setminus B}|f_k-f||\varphi|
    \le \|\varphi\|_\infty m(A)
      \sup_{A\setminus B}|f_k-f|
    \longrightarrow0.
  \]
  另一方面, Hölder 不等式给出
  \[
    \int_B|f_k-f||\varphi|
    \le \|\varphi\|_\infty\|f_k-f\|_p\,m(B)^{1/q}
    \le \|\varphi\|_\infty M\delta^{1/q}.
  \]
  令 $\delta\to0^+$, 得
  \[
    \int_E(f_k-f)\varphi\,\dif x\longrightarrow0.
  \]

  对一般的 $g\in L^q(E)$, 取有界且支集测度有限的 $\varphi$, 使
  $\|g-\varphi\|_q<\varepsilon$. 则
  \[
    \left|\int_E(f_k-f)(g-\varphi)\right|
    \le \|f_k-f\|_p\|g-\varphi\|_q
    \le M\varepsilon.
  \]
  又已知 $\int_E(f_k-f)\varphi\to0$, 所以
  \[
    \limsup_{k\to\infty}
    \left|\int_E(f_k-f)g\right|
    \le M\varepsilon.
  \]
  令 $\varepsilon\to0^+$, 即得
  \[
    \int_E f_k g\to\int_E fg. \qed
  \]
\end{solution}

% ==============================================================================
% Page 135, Exercise 15
% ==============================================================================
\section{P136 Exercise 15}

设 $f$ 和 $f_k\ (k\ge1)$ 都是 $L^p(\mathbf R)$ 中的非负函数,
$1\le p<\infty$. 此外,
\[
  \liminf_{k\to\infty} f_k(x)\ge f(x)\quad \text{a.e.},\qquad
  \|f_k\|_p\to\|f\|_p.
\]
求证:
\[
  \|f_k-f\|_p\to0.
\]

\begin{solution}
  令
  \[
    h_k(x)=\min\{f_k(x),f(x)\},\qquad
    r_k(x)=f_k(x)-h_k(x)=(f_k(x)-f(x))^+.
  \]
  由 $\liminf f_k\ge f$ a.e. 可知 $h_k(x)\to f(x)$ a.e.
  又 $0\le h_k\le f$, 故由控制收敛定理,
  \[
    \|h_k-f\|_p\to0,\qquad \|h_k\|_p^p\to\|f\|_p^p.
  \]

  因为 $f_k=h_k+r_k$ 且 $h_k,r_k\ge0$, 对 $p\ge1$ 有
  \[
    f_k^p=(h_k+r_k)^p\ge h_k^p+r_k^p.
  \]
  于是
  \[
    0\le \|r_k\|_p^p
    \le \|f_k\|_p^p-\|h_k\|_p^p\longrightarrow0.
  \]
  所以 $\|r_k\|_p\to0$. 最后
  \[
    \|f_k-f\|_p
    \le \|f_k-h_k\|_p+\|h_k-f\|_p
    =\|r_k\|_p+\|h_k-f\|_p\to0.
  \]
  证毕.
\end{solution}

\end{document}
